\section{A scalar test for the common unitary in the Bessel identity}

% BEGIN B90T1
\paragraph{B90T1. Evaluating on the native row quotient.}
Fix one selected datum in the range of R88L3. Write $E$ for its
finite-dimensional row Hilbert space, $M_u$ for the finite-dimensional
column Hilbert space of word $u$, and $M_u^\vee$ for its complex-linear
dual with the dual Hilbert norm. There are two linear maps
$T_u,S_u:E\longrightarrow M_u^\vee$ defined by
\[
 (T_uF)(\phi)=B_u(\phi,F),\qquad
 (S_uF)(\phi)=\mathrm{besS}_u(\phi,F).
\]
The native $B_u$ is complex bilinear. The second formula is also
bilinear. The row realization is conjugate-linear and enters the
conjugate-linear slot of the spectral inner product. Explicitly,
the value is $\langle s_{EE}F,P_u(\iota_{C,u}\phi)\rangle$,
with $\iota_{C,u}$ the inclusion of the ramified column space.

Use the given rational bases of the column and row spaces and their
images $(v_i)$, $(e_j)$ under the native embeddings. These images
are complex bases. R88L3 computes the coefficient matrix
$C_{u,ij}=B_u(v_i,e_j)$ by evaluating the recorded rational map
$B_{k,u}$ and applying the given field embedding. Equivalently, at
the base tuple it suffices to evaluate the period on the distinguished
lift of each $e_j$ in the good space orthogonal to the kernel of
$\iota$. R88L1--R88L3 give bijectivity onto $E$ on this complement,
and identify this lift with the one used to define $B_u$. This
requires only $\dim E$ independent lifted row representatives.

Let $H$ and $K_u$ denote the positive Hermitian Gram matrices in the
row and column bases respectively, with the inner product linear in its
second argument. The matrix of the positive row operator $T_u^*T_u$
in the row basis is
\[
                H^{-1}C_u^*(\overline{K_u})^{-1}C_u .
\]
Indeed $C_u$ maps row coordinates to the coefficients of a linear
functional in the dual column basis, whose Gram matrix is
$(\overline{K_u})^{-1}$. The same formula computes $S_u^*S_u$ with
$C_{u,ij}$ replaced by the displayed native spectral evaluations.
This uses the actual Hermitian metrics, including when the rational
bases are not orthonormal.
% END B90T1

% BEGIN B90T2
\paragraph{B90T2. Contracting the column removes its wordwise change of basis.}
For any $L:E\to M_u^\vee$ set
\[
 H(L)=
 \begin{cases}
 L^*L/\operatorname{tr}(L^*L),&L\ne0,\\
 0,&L=0.
 \end{cases}
\]
Fix a unitary $V$ of $E$. A word $u$ admits a nonzero scalar $a_u$
and a unitary $U_u$ of $M_u$ satisfying the original period identity
with this $V$ if and only if
\[
                    H(T_u)=V^*H(S_u)V.                       \tag{B90.1}
\]
Indeed, the period identity says $T_u=a_u U_u^\vee S_uV$ with
$U_u^\vee\ell=\ell\circ U_u$. Taking the row Gram and dividing by its
trace proves necessity. Conversely, equality implies that either both
maps vanish or both are nonzero, since the trace of a nonzero positive
Gram is strictly positive. In the zero case take $a_u=1$ and
$U_u=1$. Otherwise put
$c=\sqrt{\operatorname{tr}(T_u^*T_u)/\operatorname{tr}(S_u^*S_u)}$.
Then $T_u$ and $cS_uV$ have the same row Gram. The assignment
$cS_uVF\mapsto T_uF$ is therefore well defined and isometric between
their images. Extend it to a unitary of $M_u^\vee$ by matching
orthonormal bases of the two orthogonal complements. Every such
unitary is $U_u^\vee$ for a unitary of $M_u$, by finite-dimensional
duality. This gives the required identity with $a_u=c$.
% END B90T2

% BEGIN B90T3
\paragraph{B90T3. Mixed traces characterize the full native comparison.}
B87H3 established the mixed-trace criterion in orthonormal bases.
Here we apply it with the quotient lifts and the actual metrics of
B90T1. Put $d=\dim_{\mathbb C} E$ and suppose $d>0$. On the selected datum,
the native $\mathrm{BesselId}$ holds if and only if for every list of
word labels $u_1,\ldots,u_j$ of length $1\le j\le 2d^2-1$ one has
\[
 \operatorname{tr}\bigl(H(T_{u_1})\cdots H(T_{u_j})\bigr)
 =
 \operatorname{tr}\bigl(H(S_{u_1})\cdots H(S_{u_j})\bigr).    \tag{B90.2}
\]
Here labels may repeat, and their order matters. The length-one
identities require exactly that the two periods vanish for the same
labels. The formal word set is finite, so this test involves only
finitely many explicitly defined scalars for each datum.

To see sufficiency, form the pairs of row operators
$P_u=(H(T_u),H(S_u))$ in
$\operatorname{End}(E)\oplus\operatorname{End}(E)$. The spaces
spanned by products of at most $k$ such pairs, including the empty
product $(1,1)$, start in dimension one. Strict growth at each step
can occur at most $2d^2-1$ times. Equality at consecutive steps implies
stability under right multiplication by every $P_u$, so then the space
already contains all products. Thus the linear functional
$(X,Y)\mapsto\operatorname{tr} X-\operatorname{tr}Y$ vanishes on all
products by (B90.2) and vanishes on the empty product as well.

Apply the simultaneous unitary similarity criterion of Futorny,
Horn and Sergeichuk, \emph{Specht's criterion for systems of linear
mappings}, Theorem~2(a), to one vertex with one loop for each word
label. All the operators here are Hermitian, so the adjoint of a loop
contributes the very same operator. Equality of all the traces just
proved supplies one unitary $V$ satisfying (B90.1) for every label.
B90T2 recovers the wordwise unitaries and scalars. R88L3 supplies
the rational clause of $\mathrm{BesselId}$. Necessity follows from
B90T2 and invariance of trace under conjugation. If $d=0$, every
period vanishes and the identity is automatic. If $d=1$, the
length-one vanishing tests already suffice.

Consequently the arithmetic and analytic parts can be compared
directly on the native quotient, without choosing either row or column
isometries in advance. The arithmetic matrices come from the base-tuple
periods of B90T1. Establishing (B90.2) for their native spectral
counterparts on the selected family remains a period identity.
% END B90T3
